\documentclass[12pt, a4paper, oneside]{article}

\usepackage[utf8]{inputenc}
\usepackage[brazilian]{babel}

% Fonte Times New Roman (via NewTX)
\usepackage{newtxtext,newtxmath}

\usepackage{setspace}
\usepackage{geometry}
\usepackage{titlesec}
\usepackage{indentfirst}
\usepackage{booktabs}
\usepackage{array}
\usepackage{hyperref}
\usepackage{fancyhdr}

% Citações ABNT
\usepackage[alf, abnt-emphasize=bf, bibjustif, recuo=0cm]{abntex2cite}

% Configuração de página
\geometry{top=3cm, bottom=2cm, left=3cm, right=2cm}
\setlength{\parindent}{1.27cm}
\onehalfspacing

% Hiperlinks
\hypersetup{
    colorlinks=true,
    linkcolor=black,
    citecolor=black,
    urlcolor=black,
    filecolor=black,
    pdfborder={0 0 0},
    pdftitle={Pré-projeto de Mestrado},
    pdfauthor={Seu Nome},
    pdfsubject={Pré-projeto de Mestrado}
}

\begin{document}

% ============================================================
% CAPA — sem numeração
% ============================================================
\pagenumbering{gobble} % Remove número da capa

\begin{center}

    \vspace*{2cm}
    \textbf{\fontsize{12}{14}\selectfont PROGRAMA DE PÓS-GRADUAÇÃO EM INFORMÁTICA}

    \vspace{2.5cm}

    \begin{flushleft}
        \fontsize{12}{14}\selectfont
        \textbf{Nome: }

        \vspace{0.8cm}
        \textbf{CPF: }
    \end{flushleft}

    \vspace{2.5cm}

    \textbf{\fontsize{12}{14}\selectfont PROPOSTA DE PROJETO DE PESQUISA}

    \vspace{1.5cm}

    \begin{flushleft}
        \fontsize{12}{14}\selectfont
        \textbf{Título:}
    \end{flushleft}

    \vspace{1.5cm}

    \begin{flushleft}
        \fontsize{12}{14}\selectfont
        \textbf{Linha de Pesquisa:}
    \end{flushleft}

    \vspace{1.2cm}

    \begin{flushleft}
        \fontsize{12}{14}\selectfont
        \textbf{Área de Pesquisa:}
    \end{flushleft}

    \vfill

    \fontsize{12}{14}\selectfont
    \textbf{\today}

\end{center}

\newpage

% ============================================================
% INÍCIO DO TEXTO — numeração começa agora
% ============================================================

\pagenumbering{arabic} % começa contagem
\setcounter{page}{1}   % garante que começa no 1

% Cabeçalho com número no canto superior direito
\pagestyle{fancy}
\fancyhf{}
\fancyhead[R]{\thepage}
\renewcommand{\headrulewidth}{0pt}

% Formatação das seções
\titleformat{\section}
{\normalfont\fontsize{12}{14}\bfseries}{\thesection}{1em}{}
\titlespacing*{\section}{0pt}{12pt}{6pt}

\titleformat{\subsection}
{\normalfont\fontsize{12}{14}\bfseries}{\thesubsection}{1em}{}
\titlespacing*{\subsection}{0pt}{12pt}{6pt}

% ============================================================
% CONTEÚDO
% ============================================================

\section{Introdução}
Aqui deve ser apresentado o tema da pesquisa, contextualizando-o na área de estudo.

\section{Justificativa}
Argumentar sobre a importância e relevância científica do estudo.

\section{Objetivos}

\subsection{Objetivo Geral}
Descrever o propósito geral do projeto.

\subsection{Objetivos Específicos}
\begin{itemize}
    \item Identificar [aspecto 1];
    \item Avaliar [aspecto 2];
    \item Propor [aspecto 3].
\end{itemize}

\section{Revisão de Literatura}
Apresentar autores, conceitos e trabalhos relacionados \cite{Silva2019, Oliveira2020}.

\section{Metodologia}
Detalhar métodos, coleta, análise e instrumentos.

\section{Plano de Trabalho}
Descrever atividades previstas.

\section{Cronograma}

\begin{table}[h]
\centering
\fontsize{12}{14}\selectfont
\begin{tabular}{l|c|c|c|c|c|c|c}
    \toprule
    Atividades & M1 & M2 & M3 & M4 & M5 & M6 & M7--M24 \\
    \midrule
    Revisão bibliográfica & X & X & X & X &  &  &  \\
    Definição metodológica &  &  & X & X &  &  &  \\
    Coleta de dados &  &  &  &  & X & X & X \\
    Análise dos dados &  &  &  &  &  & X & X \\
    Redação da dissertação &  &  &  &  &  &  & X \\
    Revisão final e entrega &  &  &  &  &  &  & X \\
    \bottomrule
\end{tabular}
\caption{Cronograma de atividades}
\end{table}

\section{Referências}
\bibliography{referencias}

\end{document}

